% Source: https://artofproblemsolving.com/community/c3h2327153_sum_conjecture
% Stable identity: aops-topic-2327153-q01
Use partial fractions: $\dfrac{1}{k(k+2)}=\dfrac12\left(\dfrac1k-\dfrac1{k+2}\right)$. In the sum, each $-\dfrac1{k+2}$ cancels the $+\dfrac1{k+2}$ two terms later, and only the first two positive and the last two negative terms survive:
\[S_n=\dfrac12\left(1+\dfrac12-\dfrac1{n+1}-\dfrac1{n+2}\right)=\dfrac34-\dfrac{2n+3}{2(n+1)(n+2)}.\]
Over a common denominator,
\[S_n=\dfrac{3(n+1)(n+2)-2(2n+3)}{4(n+1)(n+2)}=\boxed{\dfrac{n(3n+5)}{4(n+1)(n+2)}}.\]
Check: $n=1$: $\dfrac{8}{24}=\dfrac13$; $n=2$: $\dfrac{22}{48}=\dfrac{11}{24}$; $n=3$: $\dfrac{42}{80}=\dfrac{21}{40}$; $n=4$: $\dfrac{68}{120}=\dfrac{17}{30}$.

The last value differs from the post's $\dfrac{31}{56}$: in fact $S_4=\dfrac{21}{40}+\dfrac{1}{24}=\dfrac{17}{30}$, since the fourth term is $\dfrac{1}{4\cdot6}=\dfrac1{24}$, not $\dfrac1{35}$. (The value $\dfrac{31}{56}$ comes from adding $\dfrac{1}{5\cdot7}$.)
